\section{Constructing Eligible News Observations}
\label{sec:method}

The frozen staging sample selects publishers from the English-language
\texttt{eng\_Latn} shards of Fine News. The 115-entry list includes
general-interest sources with financial coverage, and a record qualifies
if its hostname is an exact match or a dot-delimited subdomain of an
entry. We preserve the article ID, body, source URL and available
capture and publisher metadata. The recorded upstream revision and a
separate checksum-verified staging revision identify 14,807,287 distinct
IDs in 168 staged Parquet shards. The appendix records both revisions
and the separate source-comparison frames (Table~\ref{tab:app-source}).

We apply the same hostname rule retrospectively to four pinned source
frames, including FineWeb-HQ \citep{messmer2025}. In each frame we make
four independent draws with replacement, choosing a file uniformly and
then doing the same for a row group within that file. Inverse-probability
weights estimate the English and publisher-matched totals. Their ratio
gives the estimated publisher yield, with higher point estimates in the
two news collections (Table~\ref{tab:sources}). Because the sources
differ in English scope and duplicate retention, each yield describes
physical rows within that release. Estimating additional unique stories
would require matching records across releases.

\begin{table}[t]
  \centering
  \small
  \begin{tabular}{lrr}
    \toprule
    Source & Yield (\%) & 95\% interval \\
    \midrule
    Fine News & 4.50 & [2.36, 6.63] \\
    Infini-News & 6.84 & [3.37, 10.32] \\
    RefinedWeb & 0.50 & [0.16, 0.84] \\
    FineWeb-HQ & 0.69 & [0.13, 1.25] \\
    \bottomrule
  \end{tabular}
  \caption{Estimated share of physical English records that match the
    historical hostname list in each pinned source frame. English means
    the \texttt{eng\_Latn} shards for Fine News, native \texttt{eng\_Latn}
    labels for Infini-News, the declared English release for RefinedWeb,
    and native \texttt{en} labels for FineWeb-HQ. Each estimate uses
    four independent file and row-group draws with replacement.
  Intervals are exploratory and have three degrees of freedom.}
  \label{tab:sources}
\end{table}

An article is usable when it is the canonical copy and its reported
publication date is consistent with capture. We allow capture up to
seven days after that date or up to one calendar day before it to
accommodate time zones. A page dated Monday and fetched Thursday can
therefore count in the day-level corpus, while Thursday remains the
observed capture date of that version. Although Fine News retains
representatives of upstream duplicate clusters, the staged records
still contain exact-text and normalized-URL repeats. The rule uses the
earliest capture in each group. Because a repeated URL can contain
revised text, the canonical ID identifies the selected archived body.

To identify firms in the captured text, we use GLiNER2.5 Small
\citep{gliner25} to extract candidate company spans, retain their text
and offsets, and resolve the names against a dated listing and alias table.
The table combines SEC registrant identifiers and current and former
names with listing, alias and subsidiary evidence from Wikidata.
An explicit exchange ticker or Reuters code can anchor an identity and
help resolve later mentions in the same article. Alias evidence is
evaluated on the article date where validity dates are known, and the
linker retains ambiguous candidates or abstains when the evidence is
insufficient. The anchor disagreement audit includes a ticker whose
issuer in an older article about Sprint differs from a later assignment.
\citet{schwenkler2020} use Stanford CoreNLP to identify organizations
before linking firms in a news-derived network. Appendix
Table~\ref{tab:app-spans} compares GLiNER2.5 Small with Stanza, a
separate Stanford toolkit, on exact company spans in leads from the
publisher-filtered sample. FiNER-ORD also documents the difficulty of
transferring named-entity recognition to financial text \citep{shah2024}.

Because a single story can name several firms, the article--issuer
table records each accepted pair and stores the evidence for its link.
A dispute involving Gap and a Tapestry brand can concern both firms,
while an article about Tesla may mention Nvidia in passing. For the
company-focus screen we require at least three extracted mentions of
the linked issuer, with no additional relevance threshold. Each article
still counts once in the corpus and can count toward the coverage of
more than one issuer.

The pipeline also records article form, local relevance and investor-facing
sentiment for extracted mentions, and event scores for model input windows.
The overnight probe uses fixed 768-coordinate EmbeddingGemma body vectors,
whose alignment and norms we check against the frozen article IDs
\citep{embeddinggemma2025}. To examine whether a representation can
distinguish two company accounts in one story, we run a separate
pilot on 16 model-adjudicated stories with distinguishable accounts
and inputs that fit every encoder. EmbeddingGemma retrieves both
accounts correctly in 10 cases with issuer sentences and in two cases
each with a company header or a selected instruction. The pilot
motivates a larger issuer-sentence evaluation on that retrieval task
(Appendix Table~\ref{tab:app-vectors}).
